Generative AI tools were used in this work in two distinct roles. First, as
\emph{components of the research artifact itself}: large language models act as the
adversarial constructor, the user simulator, and the automated grader described in
Section~\ref{sec:methodology} and Section~\ref{sec:eval}, and every such use is
documented in the paper and the accompanying code with the specific model, version,
and prompt. Second, generative AI tools were used to assist with \emph{manuscript and
code preparation}, in four respects: prose editing and restructuring, including
condensing and reordering material we had already written; \LaTeX{} formatting,
cross-reference repair, and preparation of the alternate-format build; consistency
checking of the manuscript against our own result files, which surfaced mismatches we
then corrected; and auxiliary tooling for packaging the supplementary archives and for
normalizing in-text notation. One literature search was AI-assisted and led us to a
concurrent work that we read, verified, and cite.
The research questions and experimental design were determined by the authors.
Reported numerical results derive from the authors' experimental runs; AI assistance
also supported mechanical re-computation of stored results and preparation of tables
and the project page. The authors remain responsible for verifying the results,
figures, references, and scientific claims, and for the entire content of this paper,
in accordance with \ethicscode.
